\documentclass[./测控三级项目.tex]{subfiles} % 使用 subfiles 文档类，指定主文档

\begin{document}
	\subsection{题目重述}
		\begin{compactenum}
			\item 设计抑制题6[1]e)中的中频信号的滤波器电路，用Multisim对滤波器的频率特性进行仿真。
			\item 搭建实测电路，分别为上述低通实测电路引入一个 20mV 正弦信号，调整信号频率，在 0.1Hz-1MHz 范围内测试 30 个以上关键频率点的输出波形，记录输出信号幅值变化，用 MATLAB 描点绘制 Bode 图，用术语描述滤波器的频率特性。
		\end{compactenum}
	\subsection{题目分析与实施方案}
		\subsubsection{第一问}
			第一问要求提取题6[1]e)中的中频信号，选用了Sallen-Key结构二阶有源带通滤波器。最常见的Sallen-Key结构二阶有源带通滤波器是由一个R-C级后再连接一个C-R级构成一个带通单元，再经由R提供的正反馈得到增益单元的带通KRC滤波器。因此考虑在本题使用这个电路。 \par
			电路确定后，通过计算有源带通网络的传递函数，得到对应相频特性函数。将需要滤除的频率代入特性曲线的截止频率，便可以推算出对应电路参数，并根据实际情况，选用相关精度的标准阻值电阻。 \par
			然后，使用 Multisim 软件仿真对应电路，将题6（1）[d] 中的输出信号输入搭建的KRC有源带通滤波器上，并连接示波器。\par
		\subsubsection{第二问}
			根据第一问搭建的电路，向Sallen-Key结构二阶有源带通滤波器等对数频率，分别测试输出信号的幅值变化三十组。\par 
			使用 MATLAB 的 plot 函数，描点绘制出测试得到的参数，以得到相频伯德图曲线。\par 
			根据绘制的伯德图曲线，分析该Sallen-Key结构二阶有源带通滤波器的频率特性。
	\subsection{理论分析与仿真}
		\subsubsection{第一问}
			图\ref{fig:Sallen-Key结构二阶有源带通滤波器电路原理图}是一个Sallen-Key结构二阶有源带通滤波器电路。
			\begin{figure}[h]
				\centering
				\includegraphics[width=0.45\textwidth]{Sallen-Key结构二阶有源带通滤波器电路原理图}
				\caption{Sallen-Key结构二阶有源带通滤波器电路原理图}
				\label{fig:Sallen-Key结构二阶有源带通滤波器电路原理图}
			\end{figure}\par
			根据电路计算可得：
			\begin{align}
				\begin{dcases}
					\frac{v_a-v_P}{\frac{1}{C_2s}}=\frac{v_P}{R_2}	\\
					\frac{v_i-v_a}{R_1}+\frac{0-v_a}{\frac{1}{C_1s}}+\frac{v_o-v_a}{R_3}+\frac{0-v_P}{R_2}=0\\
					v_N=\frac{R_A}{R_A+R_B}v_o\\
					v_N=v_P
				\end{dcases}
			\end{align}\par
			整理后，可得Sallen-Key结构二阶有源带通滤波器电路的传递函数。令$s=j\omega$，并将其表示为归一化形式，可得其归一化频率$\omega_0$、归一化增益$H_0$和归一化品质因数$Q$，满足：
			\begin{gather}
				\begin{dcases}
					H_0=\frac{1+\frac{R_A}{R_B}}{1-\frac{R_AR_1}{R_BR_3}+(1+\frac{C_1}{C_2})\frac{R_1}{R_2}}\\
					\omega_0 =\sqrt{\frac{R_1+R_3}{R_1C_1R_2C_2R_3}}		\\
					Q =	\frac{\sqrt{1+\frac{R_1}{R_3}}}{(1-\frac{R_AR_1}{R_BR_3})\sqrt{\frac{R_2C_2}{R_1C_1}}+\sqrt{\frac{R_1C_2}{R_1C_1}}+\sqrt{\frac{R_1C_1}{R_2C_2}}}
				\end{dcases}
			\end{gather}\par
			为了方便参数计算，以下采用等值元件法简化。令$R_1=R_2=R_3=R$、$C_1=C_2=C$，则归一化方程组可化为：
			\begin{gather}
				\begin{dcases}
					H_0=\frac{R_A+R_B}{3R_A-R_B}	\\
					\omega_0 =\frac{\sqrt{2}}{RC}		\\
					Q =	\frac{\sqrt{2}R_A}{3R_A-R_B}
				\end{dcases}
				\label{eq:Sallen-Key结构二阶有源带通滤波器截止频率}
			\end{gather}\par
			以上理论完备，在Multisim软件中进行仿真实验。\par
			由于需要滤除的低频信号为200Hz，高频信号为1kHz，需要保留的中频信号为500Hz。因此，设置中心频率为$f_0=500$Hz，归一化品质因数$Q=0.35$。\par
			设置电容容值为$C=100$nF(Standeard Value,$5\%$)\par
			根据式\ref{eq:Sallen-Key结构二阶有源带通滤波器截止频率}，电阻阻值$R=4911.628176\Omega$，与标准电阻值表比对后，实际选用电阻阻值为$R=4.87$k$\Omega$(Standeard Value,$1\%$)的电阻。\par
			将计算得到的电阻代入得：
			\begin{gather}
				R_B=1.7655732R_A
			\end{gather}\par
			设置电阻阻值为$R_A=22.1$k$\Omega$(Standeard Value,$1\%$)，可计算得到$R_B=39.01916773$k$\Omega$，与标准电阻值表比对后，实际选用电阻阻值为$R_B=39.2$k$\Omega$(Standeard Value,$1\%$)的电阻。\par
			将要求的信号接入该电路，得到仿真电路图，如图\ref{fig:Sallen-Key结构二阶有源带通滤波器仿真原理图}所示。
			\begin{figure}[H]
				\centering
				\includegraphics[width=0.7\textwidth]{./13.第十三题示波器现象/13仿真原理图}
				\caption{Sallen-Key结构二阶有源带通滤波器仿真原理图}
				\label{fig:Sallen-Key结构二阶有源带通滤波器仿真原理图}
			\end{figure}\par
			对输入和输出信号做傅里叶变换，得到对应频谱曲线，如图\ref{fig:Sallen-Key结构二阶有源带通滤波器频谱曲线}所示。
			\begin{figure}[h]
				\centering
				\includegraphics[width=0.7\textwidth]{./13.第十三题示波器现象/13仿真频率曲线}
				\caption{Sallen-Key结构二阶有源带通滤波器频谱曲线}
				\label{fig:Sallen-Key结构二阶有源带通滤波器频谱曲线}
			\end{figure}\par
			对于低频信号（200Hz正弦信号），幅值从0.399981降到了0.244877，可求得增益为$A=20\lg(\frac{0.244877}{0.399981})=-4.261979325$dB；对于中频信号（500Hz正弦信号），幅值从0.399944升到了0.407605，可求得增益为$A=20\lg(\frac{0.407605}{0.399944})=-0.1648163383$dB；对于高频信号（1kHz正弦信号），幅值从0.0999454降到了0.0649052，可求得增益为$A=20\lg(\frac{0.0649052}{0.0999454})=-3.749666359$dB。\par
			这说明图\ref{fig:Sallen-Key结构二阶有源带通滤波器电路原理图}的带通滤波器可以对高频信号和低频信号进行滤波操作。
	\subsection{实际电路测试结果与数据分析}
		\subsubsection{第二问}
			按题目要求引入幅值为20mV的正弦信号，并调整信号频率，在0.1Hz-1MHz区间测试30个以上关键频率点，测试结果如表\ref{tab:不同频率下Sallen-Key结构二阶有源带通滤波器增益表}所示。
			\begin{table}[h!]
				\centering
				\caption{不同频率下Sallen-Key结构二阶有源带通滤波器增益表}
				\label{tab:不同频率下Sallen-Key结构二阶有源带通滤波器增益表}
				\begin{tabular}{|c|c|c|c|c|c|c|c|}
					\hline
					频率（Hz） & 增益  & 频率（Hz） & 增益 & 频率（Hz） & 增益 & 频率（Hz） & 增益\\
					\hline
					0.1&0.01&10&0.04&1000&0.81&100000&0.02\\\hline
					0.3&0.01&30&0.12&3000&0.33&300000&0.01\\\hline
					0.5&0.01&50&0.2&5000&0.2&500000&0.01\\\hline
					0.7&0.01&70&0.28&7000&0.15&700000&0.01\\\hline
					0.9&0.01&90&0.35&9000&0.12&900000&0.01\\\hline
					1&0.01&100&0.39&10000&0.11&1000000&0.01\\\hline
					3&0.02&300&0.89&30000&0.04&&\\\hline
					5&0.02&500&1.01&50000&0.03&&\\\hline
					7&0.03&700&0.96&70000&0.02&&\\\hline
					9&0.04&900&0.86&90000&0.02&&\\\hline
				\end{tabular}
			\end{table}	\par
			将表\ref{tab:不同频率下Sallen-Key结构二阶有源带通滤波器增益表}中的数据输入MATLAB中进行Bode图生成，其图像如下图\ref{fig:Sallen-Key结构二阶有源带通滤波器测试伯德图}所示。
			\begin{figure}[h]
				\centering
				\includegraphics[width=0.65\textwidth]{./13.第十三题示波器现象/Sallen-Key结构二阶有源带通滤波器测试伯德图}
				\caption{测试生成Sallen-Key结构二阶有源带通滤波器伯德图}
				\label{fig:Sallen-Key结构二阶有源带通滤波器测试伯德图}
			\end{figure}\par
			根据实际选取电阻电容值参数，可以直接得到对应传递函数。利用MATLAB软件，得到该传递函数的伯德图如图\ref{fig:Sallen-Key结构二阶有源带通滤波器实际伯德图}所示。
			\begin{figure}[H]
				\centering
				\includegraphics[width=0.65\textwidth]{./13.第十三题示波器现象/Sallen-Key结构二阶有源带通滤波器实际伯德图}
				\caption{理论生成Sallen-Key结构二阶有源带通滤波器伯德图}
				\label{fig:Sallen-Key结构二阶有源带通滤波器实际伯德图}
			\end{figure}\par
			对比图\ref{fig:Sallen-Key结构二阶有源带通滤波器测试伯德图}和图\ref{fig:Sallen-Key结构二阶有源带通滤波器实际伯德图}中的相频曲线，可以发现除了曲线前者不如后者光滑外，其余基本一致，这说明实际测试结果符合理论推断。\par
			同时也注意到，在200Hz-8kHz范围内，增益大幅变化，从0上升到1，再下降到了接近0的位置，这说明带通滤波器在该频段才具有通过信号的特点。
	\subsection{结论与收获}
		\subsubsection{结论}
		经验证，所设计Sallen-Key结构二阶有源带通滤波器基本满足题目要求，能够对中频信号进行提取，略有误差但在允许范围之内。
		\subsubsection{收获}
		\textbf{游寅龙}：\par
		理论分析层面：对电路原理图进行了查找和学习，深入理解了电路的工作原理，对该电路的理解更加深刻。\par
		动手实践层面：完成了对实测电路的数据记录和伯德图绘制，能够准确地分析数据和使用Matlab软件，增强了数据处理能力。\par
		精神层面：认识到了理论与实际是相辅相成的，更加重视电路的实际表现，增强了对实测电路调整的变通能力。\par
		\textbf{郭嘉媛}：\par
		理论分析层面：对电路相关元件进行了计算，理解了电路中各元件对滤波器参数的影响，增强了用理论知识解决问题的能力。\par
		动手实践层面：完成了对电路图的检查，能够数值电路图中各元件发挥的作用并快速准确地检查，增强了对实际电路的理解能力。\par
		精神层面：理论计算和实际电路问题检查需要逐个部分检查，培养了细心和耐心，增强了严谨的科学精神。\par
		\textbf{郑嘉宁}：\par
		理论分析层面：完成了对电路元件验算和电路的Multsim仿真，对电路原理有了更加深刻的认识也为实测结果提供了理论预期。\par
		动手实践层面：完成了对实测电路的焊接，焊接时要保证电路不出问题，需要仔细认真地焊接每个元件，培养了认真的态度和增强了焊接能力。\par
		精神层面：通过将理论预期和实测结果进行对比，体会到了理论是如何应用到实际之中，养成了理论与实际结合的工作方式。
		
	
	
\end{document}